


\begin{frame}[fragile]{Les \textit{triggers}}

Un \alert{déclencheur}  (\textit{trigger}) est  une procédure stockée qui se déclenche
 automatiquement sur certains événements.
 
 \vfill
 
\begin{redbullet}
 \item Gestion des redondances
   \item Enregistrement automatique de certains événements
  \item Gestion de contraintes complexes  (exemple: le prix d’un produit ne peut qu’augmenter)
   \item Gestion de contraintes liées à l’environnement d’exécution (restrictions sur les horaires, les utilisateurs, etc.)
\end{redbullet}
 
\vfill

\alert{Inconvénient (fort)}: les \textit{triggers} s'exécutent de manière cachée, et peuvent mener
à des cycles sans fin.
\end{frame}

\begin{frame}[fragile]{Etudions un exemple}

Un logement est constitué de chambres; sa
capacité est la somme du nombre de lits.

\vfill

On crée un \textit{trigger} qui maintient cette valeur agrégative.
 
\begin{small}
\begin{minted}{sql}
create trigger CumulCapacité
after update on Chambre
for each row
when (new.nbLits != old.nbLits)
begin
  update Logement 
  set capacité = capacité - :old.nbLits  + :new.nbLits   
  where  code = :new.codeLogement;
end;
\end{minted}
\end{small}

\end{frame}



\begin{frame}[fragile]{Les composants d'un \textit{trigger}}

Modèle basé sur la séquence Evénement-Condition-Action 
(ECA) :

\begin{redbullet}
  \item Un \textit{trigger} est déclenché par  une insertion, destruction ou modification 
   \item La condition est testée, et si elle n'est pas satisfaite, l'exécution du \textit{trigger} s'arrête
  \item L’action est effectuée  à l’aide du langage procédural  (PL/SQL)
\end{redbullet}

\vfill

\alert{De plus}\,: 
un \textit{trigger} peut manipuler 
simultanément les valeurs ancienne et nouvelle du nuplet modifié.
\end{frame}


\begin{frame}[fragile]{Les \textit{triggers} ont des inconvénients}

On peut créer des \textit{triggers} sources d'inefficacité.  Exemple

\small

\begin{small}
\begin{minted}{sql}
create trigger CumulCapaciteGlobal
after update or insert or delete on Chambre
begin
  update Logement C 
  set capacité = (select sum (nbLits) 
                       from   Chambre 
                       where  code = :new.codeLogement);
end;
\end{minted}
\end{small}
\vfill

Exécution silencieuse, donc problème difficile à détecter
\end{frame}



\begin{frame}[fragile]{Les \textit{triggers} sont dangereux\,!}

L'exécution combinée des \textit{triggers} peut avoir  des effets imprévisibles.

\vfill

La création du \textit{trigger} suivant  \alert{bloque} tout. 
\vfill

\begin{small}
\begin{minted}{sql}
create trigger MAJChambre
after update or delete or insert on Logement
begin
  update chambre
  set nbLits = (select :new.capacité / count (*)
                from  Chambre
                where codeLogement = :new.code);
end;
\end{minted}
\end{small}

\vfill

Je modifie un logement $\Rightarrow$ je modifie les chambres  $\Rightarrow$
je modifie le logement, etc. 
\end{frame}



%%%%%%%%%%%%%
\begin{frame}{À retenir}

Les \textit{triggers}\,: mécanisme puissant. Pratique, par exemple,
\vfill


\begin{redbullet}
   \item  pour contrôler ce qui se passe dans une base de données 
  \item pour mémoriser les actions sensibles (destruction)
\end{redbullet}
\vfill

Mais dangereux, à utiliser avec précaution

\begin{redbullet}
   \item L'exécution d'un \textit{trigger} est invisible\,: l'équivalent d'un
      effet de bord en programmation
   \item Un \textit{trigger}  peut être source (invisible) de lenteur
      \item La composition de plusieurs \textit{triggers} est incontrôlable
\end{redbullet}


\end{frame}
